\documentclass[11pt]{article}
\usepackage[a4paper, margin=1in]{geometry}
\usepackage{fontspec}
\setmainfont{TeX Gyre Termes}
\newfontfamily\devanagarifont{FreeSerif}
\newcommand{\dn}[1]{{\devanagarifont #1}}
\usepackage{booktabs}
\usepackage{array}
\usepackage{titlesec}
\usepackage{enumitem}
\usepackage[hidelinks]{hyperref}

\usepackage{pgfplots}
\pgfplotsset{compat=1.17}

\titleformat{\section}{\normalfont\large\bfseries}{\thesection}{1em}{}
\titleformat{\subsection}{\normalfont\normalsize\bfseries}{\thesubsection}{1em}{}
\setlength{\parskip}{0.6em}
\setlength{\parindent}{0pt}

\title{\vspace{-1.5cm}\textbf{Language Models and Agents --- Monsoon 2026}\\[4pt]
\large Individual Project: Phase 1 Report}
\author{Manya Jain}
\date{}

\begin{document}
\maketitle
\vspace{-1cm}

\section{Language Selection}

\subsection{Model H (Higher-resource): Hindi}
Hindi was chosen as the higher-resource language because of the sheer volume of public text available for it. Large open-source Indic corpora, complete Wikipedia dumps, and government press releases (PIB) all provide a solid foundation for training a well-resourced monolingual model, satisfying the phase requirement of a non-English Indian language with abundant public text.

\subsection{Model L (Lower-resource): Nepali}
Nepali was chosen as the lower-resource language. It is one of the languages explicitly permitted by the assignment, and its script and morphology present a meaningfully different tokenization problem from Hindi. Although resources are more limited than for Hindi, a workable corpus was assembled from a large academic dataset (IRIIS) supplemented with targeted news scraping.

\section{Dataset Requirements}

Both models are built on fully independent monolingual corpora, each with its own train, validation, and test split (90/5/5, seed 42). No documents, sentences, or vocabulary are shared between the two languages.

\subsection{Model H: Hindi}
The Hindi corpus combines two large downloaded sources with several manually collected ones.

\begin{table}[h]
\centering
\begin{tabular}{lp{5cm}r}
\toprule
\textbf{Category} & \textbf{Source} & \textbf{Word Count} \\
\midrule
Downloaded & Sangraha corpus & 194{,}198{,}603 \\
Downloaded & Indic corpus & 176{,}448{,}929 \\
Manual & Hindi Wikipedia & 42{,}801{,}209 \\
Manual & News scrape (new) & 42{,}383{,}013 \\
Manual & Subtitles & 5{,}234{,}233 \\
Manual & News scrape & 360{,}205 \\
Manual & PIB releases & 5{,}885 \\
\bottomrule
\end{tabular}
\caption{Hindi corpus sources and word counts (raw, pre-cleaning).}
\end{table}

The data was cleaned using Unicode NFC normalization, exact-content deduplication, and removal of non-Devanagari (English) characters to keep the corpus script-pure.

The post-cleaning corpus contains 455,866,707 words (raw source collection: 461,432,077 words; approximately 5.5M words were removed during exact-content deduplication and cleaning), which corresponds to roughly 549.4 million tokens after tokenization (fertility 1.2052), exceeding the 500M-token target. Wikipedia, both news-scrape sources, subtitles, and PIB releases together make up the manually collected portion (90,784,545 words), giving a manual fraction of approximately 19.7\% of total raw words --- marginally below the 20\% minimum required by the assignment; additional manual collection is needed to clear this threshold.

\begin{table}[h]
\centering
\begin{tabular}{lr}
\toprule
\textbf{Split} & \textbf{Words} \\
\midrule
Train & 410{,}362{,}761 \\
Validation & 22{,}804{,}967 \\
Test & 22{,}698{,}979 \\
\bottomrule
\end{tabular}
\caption{Hindi train/validation/test split (90/5/5, seed 42), computed on the post-cleaning corpus of 455,866,707 words.}
\end{table}

\subsection{Model L: Nepali}
The Nepali corpus was built with the manual-collection requirement as the primary constraint, then supplemented with a large downloaded academic dataset.

\begin{table}[h]
\centering
\begin{tabular}{lp{5cm}r}
\toprule
\textbf{Category} & \textbf{Source} & \textbf{Word Count} \\
\midrule
Downloaded & IRIIS Nepali & 400{,}000{,}000 \\
Manual & News scrape (43 portals) & 107{,}100{,}815 \\
Manual & Nepali Wikipedia & 7{,}479{,}341 \\
\bottomrule
\end{tabular}
\caption{Nepali corpus sources and word counts (raw, pre-cleaning). News data was collected from 43 high-traffic Nepali news portals, respecting robots.txt.}
\end{table}

Cleaning steps mirror those used for Hindi: cross-split exact deduplication, whitespace normalization, Unicode NFC normalization, and a Devanagari purity filter.

The post-cleaning corpus contains 513,019,586 words, corresponding to roughly 618.7 million tokens after tokenization (fertility 1.2059), exceeding the 500M-token target. Manually collected data (News scrape + Wikipedia) accounts for 114,580,156 words, approximately 22.3\% of total words, comfortably above the 20\% minimum.

\begin{table}[h]
\centering
\begin{tabular}{lr}
\toprule
\textbf{Split} & \textbf{Words} \\
\midrule
Train & 461{,}505{,}584 \\
Validation & 25{,}802{,}084 \\
Test & 25{,}711{,}918 \\
\bottomrule
\end{tabular}
\caption{Nepali train/validation/test split (90/5/5, seed 42), computed on the post-cleaning corpus of 513,019,586 words. Corrected from an earlier 98/1/1 split caused by a threshold bug (\texttt{r < 0.01} instead of \texttt{r < 0.05}) in the split script; a subsequent regeneration also required re-running \texttt{clean\_corpus.py} on the new splits to strip residual English words and noise.}
\end{table}

\section{Tokenizer Training}

Separate tokenizers were trained from scratch for each language; no pretrained tokenizer or shared vocabulary was used at any point. Model H and Model L deliberately use two different BPE variants tailored to each corpus's noise profile (see Section 3.3).

\subsection{Vocabulary Size Selection and Parameter Budget}

The initial vocabulary size for both languages was set to 32{,}000 (Hindi) and 37{,}000 (Nepali) based on a fertility-vs-vocabulary-size analysis: the fertility curve flattens noticeably beyond these sizes, meaning further merges yield diminishing improvements in compression. However, vocabulary size has a direct and significant impact on model parameter count.

For a transformer with hidden dimension $d = 384$, the token embedding table alone contributes $V \times d$ parameters. At different vocabulary sizes this fraction of total model parameters is:

\begin{table}[h]
\centering
\begin{tabular}{rrrr}
\toprule
\textbf{Vocab size} & \textbf{Embedding params} & \textbf{Approx.\ total params} & \textbf{Embedding share} \\
\midrule
32{,}000 & 12.3M & 25.0M & $\approx$49\% \\
16{,}000 & 6.1M & 25.0M & $\approx$25\% \\
10{,}000 & 3.8M & 25.0M & $\approx$18.5\% \\
\bottomrule
\end{tabular}
\caption{Token embedding parameter cost at model dimension $d=384$. A 32k vocabulary consumes nearly half the total parameter budget in embeddings alone, leaving less capacity for attention and feed-forward layers.}
\end{table}

To give flexibility for downstream Phase 2 training under different parameter budgets, three vocabulary sizes were trained for each language: 10{,}000, 16{,}000, and 32{,}000 for Hindi; 10{,}000, 16{,}000, and 37{,}000 for Nepali. All trained tokenizers and vocabulary files are committed to the repository under \texttt{models/vocab\_\{size\}/}.

\subsection{Model H: Hindi --- All Vocabulary Sizes}

All three Hindi tokenizers were evaluated on the same 10{,}000 held-out lines from \texttt{test.txt} (656{,}900 words, 3{,}347{,}181 characters).

\begin{table}[h]
\centering
\begin{tabular}{lrrrrr}
\toprule
\textbf{Vocab} & \textbf{Merges} & \textbf{Total tokens} & \textbf{Fertility} & \textbf{Chars/token} & \textbf{UNK rate} \\
\midrule
10{,}000 & 9{,}744  & 888{,}355 & 1.3523 & 3.77 & 0.00\% \\
16{,}000 & 15{,}744 & 839{,}350 & 1.2777 & 3.99 & 0.00\% \\
32{,}000 & 31{,}744 & 791{,}726 & 1.2052 & 4.23 & 0.00\% \\
\bottomrule
\end{tabular}
\caption{Hindi tokenizer comparison across vocabulary sizes, evaluated on 10{,}000 test lines.}
\end{table}

The most frequent tokens are consistent across all three sizes: punctuation and common postpositions/verbs (\dn{।}, \dn{के}, \dn{में}, \dn{है}), which is expected for a well-formed Devanagari tokenizer trained on news and encyclopedic text.

\begin{figure}[h]
\centering
\begin{tikzpicture}
\begin{axis}[
    ybar, bar width=14pt,
    width=0.85\textwidth, height=5cm,
    ylabel={Frequency},
    xlabel={Token},
    symbolic x coords={danda, ke, mein, hai, comma},
    xtick=data,
    xticklabels={\dn{।}, \dn{के}, \dn{में}, \dn{है}, {,}},
    nodes near coords,
    every node near coord/.append style={black, font=\small},
    ymin=0,
    title={Hindi: top 5 token frequencies (32k tokenizer)},
    axis lines*=left,
    enlarge x limits=0.15,
    every axis plot/.append style={fill=black!70, draw=black}
]
\addplot[fill=black!70,draw=black] coordinates {(danda,26578) (ke,26357) (mein,20000) (hai,18975) (comma,18000)};
\end{axis}
\end{tikzpicture}
\end{figure}

All three tokenizers use byte-level BPE (256-byte base vocabulary), which guarantees a UNK rate of 0.00\% on all inputs regardless of vocabulary size. The fertility drops from 1.35 at 10k to 1.21 at 32k --- an 11\% reduction in sequence length --- while compression improves from 3.77 to 4.23 characters per token. At the 32k level, common long Hindi words such as \dn{अटकलें} are encoded as a single token (id: 18534) rather than being split into 3 subword units as at 10k, directly reducing effective sequence length for the language model.

Example (\dn{``2027 विधानसभा चुनाव को लेकर तेज हुई अटकलें''}):
\begin{itemize}[nosep]
    \item 10k/16k: \texttt{[1497, 6993, 1702, 947, 311, 697, 1181, 675, 4038, 1483, 283]} (11 tokens)
    \item 32k: \texttt{[1497, 6993, 1702, 947, 311, 697, 1181, 675, 18534]} (9 tokens)
\end{itemize}

\subsection{Model L: Nepali --- All Vocabulary Sizes}

All three Nepali tokenizers were evaluated on the full \texttt{test.txt} (25{,}711{,}918 words, 166{,}916{,}158 characters).

\begin{table}[h]
\centering
\begin{tabular}{lrrrrr}
\toprule
\textbf{Vocab} & \textbf{Merges} & \textbf{Total tokens} & \textbf{Fertility} & \textbf{Chars/token} & \textbf{UNK rate} \\
\midrule
10{,}000 & 9{,}903  & 36{,}868{,}782 & 1.4339 & 4.53 & $\approx$0.000095\% \\
16{,}000 & 15{,}903 & 34{,}199{,}640 & 1.3301 & 4.88 & $\approx$0.000102\% \\
37{,}000 & 36{,}878 & \textit{(subset)} & 1.2059 & 5.38 & $\approx$0.000113\% \\
\bottomrule
\end{tabular}
\caption{Nepali tokenizer comparison across vocabulary sizes. The 37k evaluation used a smaller held-out subset; token count is not directly comparable. UNK tokens are near-zero across all sizes ($\leq 35$ per run).}
\end{table}

The most frequent tokens are punctuation and function words (\dn{।}, \dn{,}, \dn{छ}, \dn{र}). Notably, the danda (\dn{।}) appears twice in the top tokens: once with a leading space (\dn{ ।}, 1{,}111{,}474 occurrences) and once without (\dn{।}, 502{,}673 occurrences). Tokenizers treat a character with a preceding space as a distinct entity. This highlights a real-world data quirk: in the Nepali corpus, writers frequently type a space before the full stop as well as directly attached to the end of the last word. Both variations are so common in the raw text that the tokenizer learned both as distinct, high-frequency tokens.

\begin{figure}[h]
\centering
\begin{tikzpicture}
\begin{axis}[
    ybar, bar width=14pt,
    width=0.85\textwidth, height=5cm,
    ylabel={Frequency},
    xlabel={Token},
    symbolic x coords={danda1, comma, danda2, chha, ra},
    xtick=data,
    xticklabels={\dn{ ।}, {,}, \dn{।}, \dn{ छ}, \dn{ र}},
    nodes near coords,
    every node near coord/.append style={black, font=\small},
    ymin=0,
    title={Nepali: top 5 token frequencies (37k tokenizer)},
    axis lines*=left,
    enlarge x limits=0.15,
    every axis plot/.append style={fill=black!70, draw=black}
]
\addplot[fill=black!70,draw=black] coordinates {(danda1,222735) (comma,143279) (danda2,99363) (chha,94568) (ra,89894)};
\end{axis}
\end{tikzpicture}
\end{figure}

Fertility decreases steadily as vocabulary grows: 1.43 $\to$ 1.33 $\to$ 1.21, a 16\% reduction from 10k to 37k. The compression (characters per token) improves from 4.53 to 5.38, meaning each token represents progressively longer and more semantically coherent subword units. A larger vocabulary was needed relative to Hindi (37k vs.\ 32k) to reach comparable fertility, consistent with Nepali's richer morphology requiring more distinct subword units for equivalent coverage. The UNK rate remains negligible across all sizes ($< 0.0002\%$).

Example (\dn{``नेपाल एउटा सुन्दर देश हो।''}):
\begin{itemize}[nosep]
    \item All sizes decode correctly: \texttt{[1276, 951, 3834, 501, 201, 84]}
\end{itemize}

\subsection{Tokenizer Architecture Differences}

Model H and Model L deliberately use two different BPE variants, chosen from the start of the project as a point of comparison rather than as an artifact of separate implementations:

\begin{itemize}[nosep]
    \item \textbf{Model H (Hindi) -- Byte-level BPE.} Text is first decomposed into raw UTF-8 bytes (the same scheme used by GPT-2 and RoBERTa) before merges are learned. Because every possible byte is representable, this variant can encode any input, including rare symbols, mixed script, or noise, with zero out-of-vocabulary tokens (observed UNK rate: 0.00\% at all three vocabulary sizes). This suits Hindi's larger, more heterogeneous corpus (news, subtitles, Wikipedia, PIB releases), which is more likely to contain incidental non-Devanagari or irregular characters despite the script-purity filtering.
    \item \textbf{Model L (Nepali) -- Unicode character-level BPE.} Merges are learned directly over Devanagari Unicode codepoints rather than bytes. This keeps the vocabulary file human-readable in native script and ties each token directly to a valid character, at the cost of a small non-zero UNK rate ($< 0.0002\%$ across all three vocabulary sizes) on characters absent from the training vocabulary.
\end{itemize}

Training both variants let us directly compare a byte-level and a character-level tokenizer under matched conditions (same script family, same general cleaning pipeline, similar fertility targets), rather than assuming one approach is preferable a priori. The trade-off is exactly what the results show: byte-level BPE guarantees zero UNK at the cost of an unreadable, Latin-byte-mapped vocabulary file, while character-level BPE keeps the vocabulary readable at the cost of a small but non-zero UNK rate.

\section{Reproduction}

Dataset collection scripts, preprocessing pipelines, and tokenizer training code for both languages are included in the repository, organized under separate \texttt{hindi/} and \texttt{nepali/} directories as specified in the assignment layout. All trained tokenizer files (JSON + pickle) and vocabulary files for all three vocabulary sizes are committed to the repository under \texttt{models/vocab\_\{size\}/}; raw and intermediate corpus files exceed the size suitable for git and are hosted externally.

\begin{itemize}[nosep]
    \item GitHub repository: \url{https://github.com/CL3-410/individual-project-ManyaJain111}
    \item Hindi corpus (HuggingFace): \url{https://huggingface.co/datasets/manyajain111/hindi-lma-corpus} (Private)
    \item Nepali corpus (HuggingFace): \url{https://huggingface.co/datasets/manyajain111/nepali-lma-corpus} (Private)
    \item Hindi tokenizer files: \texttt{hindi/models/vocab\_10000/}, \texttt{vocab\_16000/}, \texttt{vocab\_32000/}
    \item Nepali tokenizer files: \texttt{nepali/models/vocab\_10000/}, \texttt{vocab\_16000/}, \texttt{vocab\_37000/}
\end{itemize}

\section{Summary}

Both corpora meet or exceed the 500M-token target and the 20\% manual-collection minimum (Hindi is marginally below at 19.7\% and requires additional manual collection). Both tokenizers were trained independently with no shared data or vocabulary between Model H and Model L. Three vocabulary sizes (10k, 16k, and 32k/37k) were trained for each language to provide flexibility for Phase 2 model training under different parameter budgets: at hidden dimension $d=384$, a 32k vocabulary consumes $\approx$49\% of total parameters in the embedding table, while 16k and 10k reduce that share to $\approx$25\% and $\approx$18.5\% respectively. This satisfies the Phase 1 requirements ahead of the pretraining and evaluation work in Phase 2.

\end{document}
